\chapter*{Conclusion générale}
\addcontentsline{toc}{chapter}{Conclusion générale}

Ce mémoire est parti d'une question : comment extraire de manière fiable des entités et des relations cliniques à partir de dossiers médicaux scannés de patients septiques, tout en identifiant, caractérisant et réduisant les hallucinations produites par les modèles multimodaux ? Les dossiers de la cohorte APROCCHSS, qui n'existent que sous forme d'images, en font un cas exigeant : l'information y est dispersée sur plusieurs pages, mêle texte libre, tableaux de biologie et prescriptions, et ne peut être lue qu'à travers l'image de la page.

\section*{Bilan des contributions}

La première contribution est un \textbf{référentiel d'annotation}. Le TRACE-Sepsis Annotation Guideline définit 20 types d'entités, 34 relations et les règles pour les annoter ; sa version~2.0 a été validée par le médecin référent. À partir de ce guideline, 49 dossiers (291 pages) ont été annotés à la main puis revus par le médecin, ce qui a demandé en moyenne deux heures et demie de validation par dossier. Le Gold Standard obtenu, qui conserve aussi les absences explicitement mentionnées dans les dossiers (\og pas de toux \fg{}), compte 12\,396 entités et 5\,546 relations. Il a été construit sans aucun modèle, ce qui en fait une référence indépendante des systèmes évalués. C'est ce qui manquait aux travaux antérieurs sur ce corpus, dont la qualité des extractions ne pouvait pas être mesurée finement.

La deuxième contribution est un \textbf{pipeline d'extraction multimodale} guidé par ce même guideline. Le modèle lit directement chaque page scannée, en transcrit le texte, puis en extrait les entités et les relations. Ses consignes sont générées à partir du guideline et complétées par des exemples annotés semblables à la page traitée et par des règles strictes apprises sur le Gold Standard, toujours en validation croisée. Des contrôles déterministes écartent dès l'extraction toute entité absente de la page et toute relation hors du schéma. Deux modèles, Qwen2.5-VL et Mistral Small~4, ont été comparés avec exactement le même pipeline. Mistral Small~4, retenu, atteint à la sortie de l'extraction un F1 de 67,3~\% sur les entités et de 49,7~\% sur les relations, en mode Relaxed. \acompleter{Ajouter une phrase sur l'écart avec Qwen2.5-VL une fois ses résultats disponibles.}

La troisième contribution est une \textbf{taxonomie des hallucinations} propre à l'extraction clinique multimodale. Elle distingue les informations inventées, les erreurs de type, de valeur ou d'unité, les relations mal rattachées, les incohérences structurelles et les erreurs de contexte, comme un fait nié présenté comme présent.

La quatrième contribution est \textbf{la chaîne TRACE}, qui détecte et réduit ces hallucinations après l'extraction. Ses étages ont chacun une seule responsabilité : ancrage de chaque élément dans le texte de sa page, conformité aux signatures du guideline, réparation structurelle du graphe avec une validation multi-agent, récupération des éléments oubliés, dédoublonnage, qualification des entités (unité, temporalité, négation), puis arbitrage. Le Gold Standard n'est lu que par l'étage d'évaluation, et tout ce qui en est appris l'est en validation croisée. Le mode précision de TRACE permet de choisir le compromis entre la quantité d'information gardée et sa fiabilité : en ne conservant que les entités les mieux notées, la part d'entités non justifiées par la référence descend à 8,4~\%, au prix d'une forte baisse du rappel. \acompleter{Vérifier ces valeurs avec les résultats définitifs du chapitre~6 et ajouter le résultat du mode standard.}

\section*{Réponse à la question de recherche}

Ces résultats montrent qu'une extraction fiable ne repose pas sur le seul modèle. Elle dépend d'abord d'une définition précise de ce qu'il faut extraire : le guideline sert à la fois à l'annotatrice, au modèle et aux vérifications, si bien que tous suivent les mêmes règles. Elle dépend ensuite de contrôles explicites, qui exigent que chaque élément extrait soit ancré dans le document et conforme au schéma. Elle dépend enfin d'une mesure honnête : sans Gold Standard indépendant et sans validation croisée, il serait impossible de savoir si une correction réduit réellement les hallucinations ou se contente de reproduire la référence.

Les hallucinations ne disparaissent pas pour autant. Les relations restent nettement plus difficiles que les entités, et la réduction des hallucinations se paie par une perte d'information. Plutôt qu'un système qui prétendrait tout extraire sans erreur, ce travail propose donc une chaîne dont chaque décision est traçable et dont le niveau d'exigence peut être réglé selon l'usage : exhaustivité pour l'exploration d'un dossier, précision pour la construction d'un graphe de connaissances destiné à l'aide à la décision.

\section*{Limites}

Plusieurs limites doivent être soulignées. Le Gold Standard a été annoté par une seule annotatrice, puis revu par un seul médecin : l'accord inter-annotateurs n'a pas pu être mesuré. Il ne couvre que 49 dossiers d'un même service et d'une même période, ce qui limite la généralisation des résultats. Chaque page est traitée indépendamment des autres, si bien qu'une relation dont les deux extrémités figurent sur des pages différentes ne peut pas être extraite. Le contrôle d'ancrage vérifie qu'un élément figure dans la transcription, mais pas que la transcription elle-même est juste : une valeur mal lue puis recopiée à l'identique échappe à ce contrôle. Le taux d'hallucination, défini comme la part des éléments extraits sans correspondance dans la référence, ne distingue pas à lui seul une information inventée d'une information correcte que la référence n'a pas annotée. Enfin, l'utilisation d'un modèle par une API externe impose des précautions sur la protection des données, qu'un déploiement entièrement local lèverait.

\section*{Perspectives}

Ces limites ouvrent plusieurs pistes. La première est d'élargir le Gold Standard à d'autres dossiers, avec une double annotation sur une partie du corpus pour mesurer l'accord inter-annotateurs. La deuxième est de dépasser le traitement page par page, en donnant au modèle le contexte du dossier entier pour extraire les relations entre pages et reconstituer la chronologie de la prise en charge. La troisième est d'ajouter une vérification visuelle, qui confronterait les valeurs extraites à l'image de la page et non seulement à sa transcription. La quatrième est de déployer la chaîne avec des modèles exécutés localement, sur l'infrastructure hospitalière ou universitaire. Enfin, le graphe vérifié par TRACE pourra alimenter le cadre KG-EDS et servir de base à des travaux sur le diagnostic précoce du sepsis, avec cette fois une mesure fiable de la qualité des connaissances extraites.
